\section{Moral Imagination}

\begin{quotex}
The life of the image is to be drawn from within. The life of the image is to be none other than the life of imagination. And it is of the very nature of imagination that it cannot be inculcated. There must be first of all the voluntary stirring from within. \flright{\textsc{Barfield}, \emph{Saving the Appearances}}

\end{quotex}
Although the notion of interpretations of sacred texts at levels higher than the merely literal is not new, Owen Barfield points out that there is currently a new willingness to attend to the symbolic meaning. This is shown, for example, by the very Inklings of which he was a member. The other example is \textbf{Valentin Tomberg}'s \emph{Meditations on the Tarot}. In these examples we see that the symbolic interpretation serves to deepen the understanding of the ``literal" meaning, without the intention to refute the literal meaning with the ``true" meaning.

Barfield takes this seriously enough to point out that it is actually a moral issue. Barfield explains:

\begin{quotex}
If I now maintain that these have a moral significance, and indeed a paramount moral significance, I am at once in the difficulty that the scale of values I have set up not only does not correspond with the generally accepted scale of Christian moral values, but appears to cut right across it. There are plenty of people with a natural taste for dream-psychology, or for art or literature of a symbolical nature, for sacramentalism in religion, or for other things whose meaning cannot be grasped without a movement of the imagination, who are arrogant or self-centred or in other ways no better than they should be. And conversely there are prosaic, humdrum, literal souls before whose courage and goodness we are abashed. It is not a happy task to have to maintain that, from one point of view, and that an all-important one, the former must be accounted morally superior.

\end{quotex}
This is because he regards stubborn literalness as a form of idolatry, hence a ``sin". There are the ``good" people: reliable, true, basically good, almost sinless. But Barfield asserts that the person of high imagination may be morally superior to the ``good guy". An example perhaps is William Blake who regarded imagination as a virtue. Blake was certainly a rascal in his personal life, but we remember him and not his pious next-door neighbor.

The Germans have a word ``verklemmt", which means ``being stuck", so that the outwardly virtuous may simply be inhibited or ``stuck". This is opposite to the Knight who is active in the world and nevertheless has value even if he is not so pure. Did Constantine do God's will even if he was not such a nice guy? This type goes a step further than the imaginative type, and strives to bring his imaginative vision into reality.



\flrightit{Posted on 2021-12-21 by Cologero }
